\chapter{SoC Integration}

\section{Task Description}

The task was to integrate the seven-segment display controller through APB~0
and the LED controller through AHB port~0. The unused APB~1 port was required
to return a fixed read value. The testbench was to drive the clock, release
reset, and verify that the MicroBlaze starts fetching instructions. The task
also included linting, synthesis, and implementation of the integrated design.

\section{Results}

\subsection{Controller Connections}

\texttt{basys3\_top.sv} connects both controllers to the processor subsystem
and routes their outputs to the BASYS3 display and LEDs. The base addresses
in Table~\ref{tab:address-map} are the starting addresses the processor uses
to access each peripheral.

\begin{table}[!htb]
  \centering
  \begin{tabular}{@{}lll@{}}
    \toprule
    Interface & CPU base address & Connection \\
    \midrule
    APB 0 & \texttt{0x44A00000} & Seven-segment controller \\
    AHB-Lite & \texttt{0x44A10000} & LED controller \\
    APB 1 & \texttt{0x44A20000} & Unused port with a fixed read response \\
    \bottomrule
  \end{tabular}
  \caption{Peripheral addresses in the Lab 4 block design.}
  \label{tab:address-map}
\end{table}

The archived RTL returns \texttt{0xA5A5A5A5} on APB~1. The Lab 4 handout
specifies \texttt{0xA5A55A5A} for this response.

\subsection{Processor Startup Simulation}

The archived \texttt{basys3\_tb} run released reset at 100\,ns and reached
the end of the test at 100.1\,$\mu$s. The processor startup check requires
completed instruction reads on the local memory bus after reset release.

\evidencefigure{lab4-startup-waveform.png}{MicroBlaze instruction reads after reset release in the Lab 4 simulation.}{fig:lab4-startup}{Capture reset release and the first completed instruction reads in the \texttt{basys3\_tb} waveform.}

\subsection{Synthesis and Implementation}

The archived Lab 4 Vivado run reports 2,383 slice LUTs, 2,229 flip-flops,
and 32 block RAM tiles for \texttt{basys3\_top}. All 3,868 routable nets
were routed, with zero routing errors.
